%% CPP 2027 Submission — Empirical Foundations of Pat:
%% transformer OOD observations, structural laws, and machine-checked theorems
\documentclass[sigplan,10pt,anonymous,review]{acmart}
\settopmatter{printfolios=true,printacmref=false}

\AtBeginDocument{%
  \providecommand\BibTeX{{\normalfont B\scshape ib\TeX}}}

\settopmatter{printacmref=false}
\renewcommand\footnotetextcopyrightpermission[1]{}
\pagestyle{plain}

\usepackage{microtype}
\usepackage{booktabs}
\usepackage{amsmath}
% XeTeX (tectonic): acmart loads newtxmath/libertine which already defines \Bbbk;
% reset it so amssymb's definition is accepted.
\let\Bbbk\relax
\usepackage{amssymb}

\begin{document}

\title{Empirical Foundations of Pat: Transformer OOD Observations,\\
Structural Laws, and Machine-Checked Continuum Theorems}

\author{Anonymous}
\affiliation{%
  \institution{Anonymous}
}
\email{anonymous@example.org}

\begin{abstract}
We record the empirical route by which the Pat framework is grounded:
starting from controlled transformer training observations, extracting
structural laws, and ending in machine-checked theorems. The starting
observation is a two-pole out-of-distribution (OOD) dichotomy, reproduced
in 45 training runs: a low-epoch transformer reaches training accuracy
$1.000$ on successor-notation arithmetic, yet OOD accuracy is $32/32$ on
length extrapolation in the trained notation and $0/32$ on the same
structure written in an untrained notation of the same basepoint; adding
explicit equivalence statements bridges the two notations ($0/32 \to
32/32$). From this and supporting experiments we extract five structural
laws (anchor binding, notation translatability, shift invariance,
structure-bearing representation, definition-layer invisibility). The
laws map, as an observational correspondence, to the Pat framework's
paired-direction declaration, phase-relation locking, and phase-magnitude
interlock. A completeness criterion --- the decoupling operator $D$
(exclusion + cancellation + layering) implying completeness at any
order, element, and subject-object --- organizes the correspondence, and
three mapping judgments (symbol-norm, presentation-legality,
intuition-precision) give the layer-to-layer verdicts that the OOD
two-pole dichotomy exhibits. The framework chain is formalized in Lean
4 / mathlib: twenty-two theorems with zero errors and zero
\texttt{sorry}, from the interlock matrix (determinant
$\theta \cdot (r + 1/r) \neq 0$) through the four-phase interlock, the
countable pat grid, and the density theorem, to the continuum closure
theorems (the continuum $[0, 2\pi]$ is the closure of the pat grid; an
endogenous version constructs the same closure from two paired
basepoints and midpoint reduction). The honest boundary is stated
exactly: the empirical-to-framework mapping is an observational
correspondence, not a deductive implication, and the training results
are reproducible from the artifact (one command, Python + PyTorch).
\end{abstract}

\keywords{Lean, formalization, transformer, out-of-distribution, pat grid, continuum, empirical}

\maketitle

\subsection*{Verification map}

\begin{center}
\small
\begin{tabular}{@{}p{0.22\textwidth}p{0.72\textwidth}@{}}
\toprule
\textbf{Main theorems} &
Continuum closure: every $x \in [0, 2\pi]$ lies in the closure of the
pat grid (and in the closure of the endogenous constructible grid);
interlock: the phase-magnitude matrix is nonsingular; mapping
judgments: resolvability verdicts. \\
\midrule
\textbf{Assumptions} &
Lean 4.32.2 + mathlib v4.32.2; standard Lean axioms
(\texttt{propext}, \texttt{Classical.choice}, \texttt{Quot.sound}
per the mathlib dependency); no additional axioms, no
\texttt{sorry}. \\
\midrule
\textbf{Core lemma chain} &
R143 interlock ($\det \neq 0$) $\to$ R149 four-phase interlock $\to$
R150 countable grid + density $\to$ R151 continuum closure $\to$
R163 endogenous closure (two basepoints + midpoint reduction). \\
\midrule
\textbf{Lean theorems} &
\texttt{MutualLocking.lean} (3), \texttt{Pat4Phase.lean} (4),
\texttt{PatCountableInfinitPhaseUnification.lean} (5),
\texttt{ContinuumPatGrid.lean} (1),
\texttt{ContinuumConstructible.lean} (4),
\texttt{MappingJudgmentTheorems.lean} (5) --- 22 total. \\
\midrule
\textbf{Repository / version} &
artifact in this submission; \texttt{lake build} passes; mathlib
v4.32.2 pinned by \texttt{lean-toolchain}; 26-file minimal dependency
closure. \\
\bottomrule
\end{tabular}
\end{center}

\section{Introduction}

The Pat framework is a family of claims about direction declaration,
phase locking, and symmetric-pair reduction, developed in a program of
small transformer experiments. Small transformers trained on
algorithmic datasets exhibit rich and counterintuitive generalization
behavior: delayed generalization (grokking) on small algorithmic
datasets~\cite{power2022grokking}, systematic failures on
out-of-distribution inputs~\cite{deletang2022chomsky,
madsen2020arithmetic}, and length extrapolation that depends on the
position encoding~\cite{press2021alibi, su2021roformer,
duan2023interpolation}; in-context arithmetic emerges at
scale~\cite{brown2020gpt3}. This paper records the empirical route
that grounds the framework: what was observed, what structural laws
were extracted, and which claims were machine-checked. The route is a
chain of four links.

\paragraph{Empirical starting point.} A low-epoch transformer trained on
successor-notation arithmetic (addition up to 9, multiplication up to 4,
fully enumerated with true and false labels) reaches training accuracy
$1.000$ in every one of 45 runs (5 groups $\times$ 3 seeds $\times$ 3
independent executions). Yet out-of-distribution (OOD) accuracy
dichotomizes: $32/32$ on length extrapolation in the trained notation,
$0/32$ on the same structure written in an untrained second notation of
the same basepoint, and $32/32$ after explicit equivalence statements
bridge the two notations. Training completeness does not imply
extrapolation completeness; this is the starting observation from which
everything else follows.

\paragraph{Structural laws.} From this and the supporting experiments
we extract five laws about what training intuition is and is not:
operations bind to a basepoint (anchor binding); notations of the same
anchor are translatable when equivalence is declared (notation
translatability); one step of successor extrapolates along the chain
(shift invariance); structure must be fully expanded in the sequence
(structure-bearing representation); and the definition layer of tokens
is invisible to training (definition-layer invisibility). The meta-law
is: training intuition is representation-bound statistics.

\paragraph{Framework correspondence.} The laws correspond, as an
observational correspondence rather than a deductive implication, to
framework claims: anchor binding to paired basepoint declaration
(R136), notation translatability to symmetric-pair reduction (R143),
and the collapse of undeclared notations to the self-referential
degenerate case (R138). A completeness criterion organizes the
correspondence: the decoupling operator $D$ --- exclusion, cancellation,
layering --- implies completeness at any order, any element, and any
subject-object, exhibited as perfect OOD behavior; and three mapping
judgments (symbol-norm, presentation-legality, intuition-precision)
give the layer-to-layer verdicts that the OOD two-pole dichotomy
exhibits.

\paragraph{Machine-checked theorems.} The framework chain is formalized
in Lean 4 / mathlib: twenty-two theorems, zero errors, zero
\texttt{sorry}, from the interlock matrix to the continuum closure
theorems. The theorem chain is: the interlock matrix is nonsingular
(det $= \theta \cdot (r + 1/r) \neq 0$); the four phases
$\{1, -1, i, -i\}$ pairwise interlock; every phase in
$[0, 2\pi]$ is uniformly approximated by the countable pat grid
$\{2\pi j / N\}$; the continuum is the closure of the pat grid; and an
endogenous construction from two paired basepoints and midpoint
reduction reaches the same closure.

The paper is organized as the route itself: Section~2 reports the
empirical starting point; Section~3 states the structural laws;
Section~4 gives the framework correspondence; Section~5 states the
completeness criterion and the mapping judgments; Section~6 presents
the machine-checked theorems; Section~7 describes the exercise
series; Section~8 states the honest boundary; Section~9 concludes.

\section{The empirical starting point}

\subsection{Setup and judgment discipline}

All experiments use the same training runtime (included in the
artifact), a small transformer~\cite{vaswani2017attention}
(dimension 64, 2 layers, 60 epochs,
learning rate $10^{-3}$, batch 512) over tokenized arithmetic
sequences. Tokens are fully expanded successor chains (true numerals):
the chain length carries the number. The judgment discipline is strict:
a sample is correct only if the full sequence is reconstructed token by
token (end-truth judgment); positional accuracy is not trusted, because
aggregate accuracy masks collapse points. Every reported count is a
full-sequence count.

\subsection{Observation 1: training accuracy 1.000, OOD dichotomizes}

Training completes perfectly in all 45 runs (5 groups $\times$ 3 seeds
$\times$ 3 executions) --- no exception. OOD, however, splits into two
poles: some probes score $32/32$, others $0/32$. Training completeness
$\neq$ extrapolation completeness. This dichotomy is the fixed point
from which all subsequent observations are made.

\subsection{Observation 2: same-notation extrapolation is 32/32}

Group \textbf{1a}: train on the notation $\mathit{succ}_a$ (addition
$\le 9$, multiplication $\le 4$), probe OOD by length extrapolation in
$\mathit{succ}_a$ (addition $10$--$12$, multiplication $5$--$6$):
$32/32$. The chain-length rule (chain concatenation) holds beyond the
training coverage. Within a representation, statistics extrapolate;
this is consistent with prior reports of length generalization in
arithmetic transformers~\cite{duan2023interpolation,
madsen2020arithmetic}.

\subsection{Observation 3: same-basepoint zero-shot is 0/32}

Group \textbf{1b}: the training samples are identical to 1a; only the
OOD probe changes notation, from $\mathit{succ}_a$ to $\mathit{succ}_b$
--- the same successor structure, same basepoint, a different spelling.
The probe scores $0/32$: complete collapse. $\mathit{succ}_b$ is
``the same successor written differently''; zero-shot does not
transfer. Training intuition binds to the notation, not to the
semantics: the meaning of \emph{successor} never entered the training.

\subsection{Observation 4: equivalence statements bridge, 0/32 to 32/32}

Group \textbf{2}: add equivalence contrast samples to the training set
($\mathit{succ}_a$-chain $n$ = $\mathit{succ}_b$-chain $n$, $n =
0..9$) plus anchored wrong samples (shown by a control group to be
irrelevant). The OOD probe on $\mathit{succ}_b$ now scores $32/32$.
Explicit statements that the two notations are equivalent bridge the
notations: $0/32 \to 32/32$. An equivalence statement is a paired
declaration; it is the bridge tool.

\subsection{Reproducibility}

All three executions (two in the original environment, one inside the
self-contained package) reproduce the same per-seed counts: $32/32$,
$0/32$, $32/32$ for groups 1a, 1b, 2 respectively, with the control
groups 2' (bridge does not damage the original notation) and 4
(symmetry-special training) also $32/32$. The core two-pole contrast
($0/32 \leftrightarrow 32/32$) is stable across seeds and executions.

\section{Five structural laws}

\subsection{Law 1: anchor binding (E11)}

Shifting the basepoint (anchor drift) collapses zero-shot to $0/32$;
each anchor requires its own training; explicit equivalence samples
between anchors do not transfer the operation rules. Anchor type makes
no difference (a value-zero anchor and a basepoint anchor drift
identically). Every reference frame's operations must be established
independently.

\subsection{Law 2: notation translatability (E12)}

Different notations of the same anchor are bridgeable by equivalence
statements ($0/32 \to 32/32$); the same notation across anchors is not
bridgeable (E11, $0/32$). The anchor (basepoint) is the structural
identity to which operations bind; a notation is a translatable
representation within an anchor. Contrast samples (equivalence
statements) are the bridge tool.

\subsection{Law 3: shift invariance (E9)}

Training only $\mathit{succ}(0) = 1$ (two samples) extrapolates to
$\mathit{succ}(1..7)$: $14/14$. The step $0 \to 1$ and the step
$1 \to 2$ are the same successor; this is the mechanism of chain-length
extrapolation.

\subsection{Law 4: structure-bearing representation (E6)}

Position-value bare symbols (implicit place weight, nominal position
tokens) collapse ($0/30$); true numerals (fully expanded successor
chains, chain length = value) are stable ($30/30$). Structure must be
fully expanded in the sequence; it cannot be left to intuition.

\subsection{Law 5: definition-layer invisibility (E4/E5)}

Four definitions of 0 (postulate, iteration zero, empty-chain default,
cancellation law) and two designs of successor (explicit concept,
atomic counting) train equivalently. Training intuition is sensitive
only to the representation form (sequence structure); the token
definition layer (where the semantics comes from) is invisible.

\subsection{The meta-law: training intuition is representation-bound statistics}

Strip one representation and probe the same value in another: total
collapse ($0/32$). Joint training on several representations does not
induce cross-representation invariants (a third representation still
scores $0/32$). An OOD score of $1.0$ proves only that statistics
extrapolate within that representation; it says nothing about
structure. Structure (value, semantics) is the object of the
formalization layer (Lean, definitions); the training layer cannot
reach it.

\section{The framework correspondence}

The five laws correspond to framework claims. We state the
correspondence explicitly and mark it as observational: each framework
concept has a statistical manifestation, but the manifestation does not
imply the concept.

\paragraph{Anchor binding $\leftrightarrow$ paired basepoint declaration (R136).}
Empirically, operations bind to an anchor: shifting it collapses
($0/32$), changing notation without declaration collapses ($0/32$).
In the framework, a basepoint cannot be used alone; directions must be
declared in symmetric pairs at once (a single direction is a privilege
pollution). The empirical $0/32$ is the statistical manifestation of
an undeclared pair: training $\mathit{succ}_a$ without declaring
$\mathit{succ}_b$ leaves the $\mathit{succ}_b$ side without an anchor,
and it collapses.

\paragraph{Equivalence statements $\leftrightarrow$ symmetric-pair reduction (R143).}
Empirically, equivalence contrast samples bridge two notations
($0/32 \to 32/32$). In the framework, a symmetric pair
$\{\mathit{succ}_a, \mathit{succ}_b\}$ reduces to 1: the numerical pair
$r \cdot (1/r) = 1$ and the phase pair $e^{i\theta} \cdot e^{-i\theta}
= 1$ are the two decompositions of the unit. The bridge is the
statistical manifestation of symmetric-pair reduction: once declared,
$\mathit{succ}_b$ is solvable from $\mathit{succ}_a$ and back.

\paragraph{Interlock: bidirectional solvability (R143).}
After bridging, both directions work (groups 2 and 2', both $32/32$).
In the framework, interlock is equivalent to nonsingularity of the
interlock matrix ($\det = \theta \cdot (r + 1/r) \neq 0$): from either
coordinate (phase or magnitude) the other is solvable. Bidirectional
translatability is the statistical manifestation of interlock
nonsingularity.

\paragraph{Undeclared = collapse (R138).}
An undeclared equivalence collapses to $0/32$: information is
unreachable across representations. In the framework, an unlocked phase
relation is the self-referential degenerate case: every operation on it
returns itself (absorption). The $0/32$ is the statistical
manifestation of the degenerate collapse: the undeclared side's
information is unreachable, and the loop closes on itself.

\paragraph{Shift invariance $\leftrightarrow$ successor identity (R137).}
Empirically $\mathit{succ}(0)$ extrapolates to $\mathit{succ}(1..7)$.
In the framework, a declared direction yields a non-collapsing chain:
each successor step is the same direction. Chain-length extrapolation
is the statistical manifestation of chain injectivity.

\section{The completeness criterion and the mapping judgments}

\subsection{The completeness criterion}

The correspondence above is organized by a completeness criterion: the
decoupling operator $D$ --- exclusion (remove what does not belong),
cancellation (pair off and cancel what reduces), layering (separate
what plays different roles) --- implies completeness at any order, any
element, and any subject-object, exhibited as perfect OOD behavior.
The empirical program verifies the criterion's clauses: at the order
level, algebraic axioms $12/12$, involutions $12/12$, truth/false
separation $6/6$ (token-level, per the reporting discipline); at the
element level, the orthogonality-declaration bridging succeeds
same-anchor ($32/32$) and a double declaration (orthogonality +
equivalence) restores stability across seeds ($5/5$); at the
subject-object level, direction localization with two symmetry groups
scores $16/16 \times 3$ seeds, and a single declaration is
insufficient --- two groups of symmetries are necessary to localize a
direction (Cartan--Dieudonn\'e: rotations are composites of two
reflections; a single reflection only defines $\pm$).

\subsection{The three mapping judgments}

The OOD two-pole dichotomy is the training-layer exhibition of three
layer-to-layer judgments about the token system (five layers:
basepostulate B, concept C, symbol S, grammar G, presentation P).
Each judgment decides, for a multi-to-multi mapping between layers,
whether every target is resolvable.

\paragraph{Symbol-norm judgment (S $\to$ C/I).}
How is the multi-to-multi mapping from symbols to logic and intuition
judged? Every target must be resolvable: ambiguity converges by
context (no context = all candidates; given context = one);
unresolvable ambiguity or collision is non-norm. The two symmetric
directions are symbol $\to$ logic and symbol $\to$ intuition. The
experimental exhibition is the structure-vs-symbol contrast: mapped
symbols $30/30$, bare symbols $0/30$ (E6).

\paragraph{Presentation-legality judgment (C/I $\to$ P).}
How is the multi-to-multi mapping from logic and intuition to
presentation judged? Every presentation source must be resolvable:
multiple presentations of one source are bridgeable by equivalence
statements; numeric coincidence without explicit declaration is
illegal (the supervision signal carries no direction information).
The experimental exhibitions are the two poles of the dichotomy:
undeclared notation collapse ($0/32$) and equivalence bridging
($32/32$).

\paragraph{Intuition-precision judgment (I $\to$ C/P).}
How is the multi-to-multi mapping from intuition to logic and
presentation judged? The intuition-token interface must be
definitionally resolvable (it must not impersonate a logic concept);
intuition-path precision is judged by the two-pole comparison
(in-representation OOD $1.0$, cross-representation $0$); polarity
confusion events are the canonical failure. The experimental
exhibition is exactly the two-pole dichotomy with the bridge as the
resolvability restoration ($0/32 \to 32/32$).

A fourth judgment (symbol $\to$ presentation: notation context
precedence and associativity resolve; glyph identification is unique;
round trips are lossless) completes the five-layer coverage. All four
judgments share one structure: resolvability of every target of a
multi-to-multi mapping, judged by declaration, context, or two-pole
comparison.

\section{Machine-checked theorems}

The theorems are formalized in Lean 4~\cite{lean4} against
mathlib~\cite{mathlib2020, mathlib2026}. Recent systems couple language
models with proof assistants for formal mathematics --- from
curriculum learning over formal statements~\cite{polu2022curriculum,
han2021gptf} and the MiniF2F benchmark~\cite{zheng2021minif2f}, to
retrieval-augmented and copilot-style proving~\cite{yang2023leandojo,
song2024leancopilot}, interleaved thinking and
proving~\cite{lin2024leantar}, large-scale synthetic
data~\cite{xin2024deepseekprover, azerbayev2023llemma}, and guiding
provers with informal proofs~\cite{jiang2022draftsketch}. The present
development is a small, fully machine-checked library produced in this
style (Section~8).

\subsection{The interlock matrix (R143, MutualLocking.lean)}

A declaration is two groups of symmetries: a phase (direction) pair
$(\theta, -\theta)$ and a magnitude (distance) pair $(r, 1/r)$. The
two paired vectors are a $2 \times 2$ matrix
\[
  \begin{pmatrix} \theta & r \\ -\theta & 1/r \end{pmatrix},
  \qquad \det = \theta \cdot (r + 1/r).
\]
\verb|mutual_lock_invertible|: for $\theta \neq 0$ and $r > 0$ the
determinant is nonzero --- the matrix is nonsingular, so either
coordinate is solvable from the other. Phase and magnitude are mutually
locked; no third layer of locking is needed. Supporting theorems:
\verb|magnitude_locks_phase_round_trip| (the normalized displacement
$(pat0 + d - pat0)/\|pat0 + d - pat0\| = d$ recovers the direction
from the locked magnitude position) and \verb|magnitude_pair_log_mirror|
($\log(1/r) = -\log r$: the magnitude pair is the log mirror pair).

\subsection{Four-phase interlock (R149, Pat4Phase.lean)}

The unit's two symmetry decompositions combine: the numerical pair
$a \cdot (1/a) = 1$ and the phase pair $e^{i\theta} \cdot e^{-i\theta}
= 1$. \verb|quadriphase_interlock|: the four phases $\{1, -1, i, -i\}$
pairwise interlock (magnitude pair + phase pair), normalizing to
$S^3$. Supporting theorems establish axis orthogonality
(\verb|axis_pair_orthogonal|) and the extrapolation to the pat circle
(\verb|extrapolation_to_pat_circle|, \verb|infinite_isomorphic_extrapolation|).

\subsection{Countable grid and density (R150, PatCountableInfinitPhaseUnification.lean)}

The pat grid is the union of all $N$-slot rings:
\[
  \mathrm{patGrid} = \big\{\, 2\pi j / N \;\big|\; 0 < N,\; j \le N
  \,\big\}.
\]
\verb|pat_grid_countable|: the grid is countable (it is the image of
$\mathbb{N} \times \mathbb{N}$). \verb|pat_phase_unification|: every
phase $\theta \in [0, 2\pi]$ is approximated to arbitrary precision by
a grid point ($\forall \varepsilon > 0,\ \exists x \in
\mathrm{patGrid},\ |\theta - x| \le \varepsilon$); the quantization
error is $\le \pi/N$, so $N \ge \pi/\varepsilon$ suffices. This is the
bridge theorem from the finite construction to the continuum.

\subsection{Continuum closure (R151, ContinuumPatGrid.lean)}

\verb|continuum_in_pat_grid_closure|: every $x \in [0, 2\pi]$ belongs
to the closure of the pat grid. The proof uses the density theorem and
the closure criterion directly (no sequence or limit lemmas). The
continuum is not a clean object; it is the closure of the countable
reachable grid --- a generated structure's limit.

\subsection{Endogenous closure (R163, ContinuumConstructible.lean)}

The closure is re-derived from endogenous generation: two paired
basepoints ($base_0 = 0$, $base_1 = 2\pi$) and the symmetry reduction
rule (the reduction point of a symmetric pair is its midpoint)
generate the grid
\[
  \mathrm{patGrid}' = \big\{\, 2\pi \cdot k / 2^m \;\big|\; 0 \le k
  \le 2^m \,\big\}
\]
(dyadic rationals $\times$ $2\pi$). \verb|dyadic_InPat|: every dyadic
point is generated; \verb|patGrid'_dense|: the generated grid is dense
in $[0, 2\pi]$; \verb|continuum_in_constructible_closure|: every $x
\in [0, 2\pi]$ lies in the closure of the generated grid. The
definition layer uses no natural numbers: generation is closure under
midpoint from the two paired basepoints, and the dyadic representation
is a representation theorem, not an input. The construction mirrors
the paired-declaration discipline: two basepoints, declared as a pair;
symmetric pairs reduce to their midpoint.

\subsection{Mapping judgments (MappingJudgmentTheorems.lean)}

Five theorems formalize the resolvability verdicts of Section~5:
\verb|context_convergence_resolvable| (context convergence implies the
existence of a resolution verdict for the symbol mapping),
\verb|equivalence_bridge_judgment| (an equivalence declaration yields a
bridge verdict), \verb|conflict_declared_resolvable| (a declared
conflict yields an injective resolvability verdict),
\verb|round_trip_precise_injective| (lossless round trip $B \circ A =
\mathrm{id}$ implies the assembly map is injective --- symbols are
resolvable), and \verb|polar_judgment_separated| (the two-pole
judgment separates: in-representation truth and cross-representation
falsity are judged by the same verdict structure). All five compile
with zero errors and zero \texttt{sorry}.

\section{The exercise series}

The route is accompanied by an independent verification face: the
homework-exercise series --- twenty-five observation reports
(Exercises I--XXII, the Pat0 exercise, and the series overview), each
with its own Lean formalization, all zero \texttt{sorry}. The method
of the series is the ``0pat'' discipline: take a pat observation,
walk the genealogy map from the pat concept to its standard-
mathematics coordinate, formalize the equivalence along the way, and
keep the pat vocabulary out of the proof. Each exercise is a small
independent formalization; together they form a corpus of
machine-checked statements that the framework's concepts reduce to.
The full series (reports, TeX sources, Lean code, and overview) is
included in the artifact under \texttt{exercises/}. The series also
records the oldest empirical anchor of the route: the first
generalization observation, dated 2026-07-29 in the session log.

\section{Honest boundary}

The boundary is stated exactly, in four points.

\paragraph{The empirical-to-framework mapping is observational.}
The correspondences of Section~4 (statistical manifestations) are not
deductive implications: no theorem states that training behavior
implies a framework claim. The Lean theorems hold independently of the
training observations; the observations provide motivation and
correspondence.

\paragraph{No structural claim about the model.}
Training intuition is representation-bound statistics. The paper does
not claim the transformer learned structure; it reports what the
statistics do and do not transfer to.

\paragraph{The proofs were developed in a human--AI collaboration.}
The direction, the framework claims, the observation judgments, and
the corrections are human-authored; the Lean proof scripts were
drafted by an AI assistant and checked by the machine. The
collaboration is part of the methodology, not of the proofs: every
theorem stands on the machine check alone, and no claim depends on
the drafting process. This mirrors the growing practice of
AI-assisted formalization~\cite{yang2023leandojo, song2024leancopilot,
lin2024leantar, xin2024deepseekprover}.

\paragraph{Novelty discipline.}
The classical components are marked KNOWN: the interlock matrix is a
2--2 nonsingularity statement over the reals; the density of dyadic
rationals in the interval is classical; the closure theorems are
reformulations of standard density results in the framework's
vocabulary. The claims R136--R163 carry \texttt{novelty\_status}
annotations (KNOWN or NOVELTY\_UNASSESSED) in the artifact's claim
ledger. The empirical observations are reported as observations, not
as theorems.

\section{Conclusion}

The route from transformer empiricism to the Pat framework and the
continuum is complete and recorded: the two-pole OOD dichotomy
($0/32 \leftrightarrow 32/32$, 45 runs reproduced) yields five
structural laws (anchor binding, notation translatability, shift
invariance, structure-bearing representation, definition-layer
invisibility); the laws correspond observationally to the framework's
paired declaration, phase locking, and interlock; a completeness
criterion (the decoupling operator $D$) organizes the correspondence;
and the framework chain is machine-checked in Lean 4 / mathlib ---
twenty-two theorems, zero \texttt{sorry}, from the interlock matrix
through the countable pat grid and its density to the continuum
closure theorems.

In one sentence: the training layer says ``representation-bound
statistics,'' the formalization layer says ``paired-declaration
structure,'' and the two meet at the continuum closure: a
representation is an anchor's notation, and the continuum is the
anchor's closure.

\begin{acks}
We thank the Lean and mathlib communities for the toolchain.
\end{acks}

\bibliographystyle{ACM-Reference-Format}
\bibliography{refs}

\end{document}
